%% ch2_background.tex — Chapter 2: Background and Related Work
%% HSP-Agile Thesis, Ye Xujun, ECNU, July 2026

\chapter{Background and Related Work}
\label{ch:background}

This chapter surveys the technical foundations and related prior work
that underpin the HSP-Agile pipeline.
We begin with the formal language framework in which our work is
situated---SOFL and its agile variant---then describe the Functional
Specification Format (FSF) at the centre of the pipeline, followed by
reviews of LLM-based code generation, counterexample-guided synthesis,
SMT solving, requirement-level pattern detection, and mutation testing.
The chapter closes with a structured comparison of related approaches
and identifies the research gap that HSP-Agile addresses.

% ─────────────────────────────────────────────────────────────────────────────
\section{SOFL and Agile-SOFL}
\label{sec:bg:sofl}
% ─────────────────────────────────────────────────────────────────────────────

\subsection{The SOFL Formal Method}

Structured Object-Oriented Formal Language (SOFL) was introduced by
\citet{liu1997sofl} as a \emph{process-oriented} formal method that
synthesises three influential formalisms: Z notation for specification
of abstract data types and state transitions, data flow diagrams for
modelling process decomposition and information flow, and object-oriented
concepts for structuring large systems into modular, reusable components.
The motivation behind this synthesis was practical: pure mathematical
formalisms such as Z or VDM provide strong logical foundations but are
perceived as inaccessible to practising engineers accustomed to diagrammatic
notations, while purely graphical methods lack the semantic precision
required for verification.
SOFL bridges this gap by providing a uniform textual language whose
constructs map directly to well-understood notations while supporting
mechanised analysis.

The cornerstone of SOFL's development paradigm is \emph{stepwise refinement}
\citep{back1998refinement}.
A software system is described at successively more concrete levels: from an
abstract specification that captures \emph{what} a function must do, through
an intermediate design that begins to constrain \emph{how} it does so, to a
concrete implementation.
At each level the developer must demonstrate that the refinement step
preserves the properties of the more abstract specification, a discipline
that prevents the accumulation of specification drift over the development
lifecycle.
SOFL has been applied to a range of industrial domains including railway
interlocking systems, medical device firmware, and embedded control software,
with \citet{liu2014sofl} providing a comprehensive textbook account of the
method alongside case studies.

A key element of SOFL semantics is the \emph{process specification}, which
describes a computation by enumerating named scenarios, each governed by a
guard condition that delimits its domain of applicability and a body that
defines outputs in terms of inputs.
This scenario-based structure is the direct ancestor of the Functional
Specification Format described in Section~\ref{sec:bg:fsf}, and it is the
primary unit around which the HSP-Agile pipeline is organised.

\subsection{Agile-SOFL: Lightweight Formal Methods in Sprints}

Classical applications of SOFL follow a plan-driven development model in which
the full specification is completed before any implementation work begins.
While this discipline is appropriate for safety-critical systems with well-understood
requirements, it is at odds with the \emph{incremental, iterative} delivery
cycle of Agile methods \citep{beck2001xp}, in which requirements are
elaborated progressively and working software is delivered at the end of each
sprint.

\emph{Agile-SOFL} \citep{li2022agilesofl} was developed to reconcile formal
specification with agile practice.
The central insight is that formal rigour need not be applied uniformly across an
entire system from the outset; instead, it can be concentrated on the
\emph{behavioural contracts} of individual functions that are part of the current
sprint's scope.
Each sprint produces a set of FSF specifications alongside executable code;
the pipeline then automatically checks conformance, generates test cases,
and flags violations before the sprint review.
This confines the cost of formal analysis to a bounded, sprint-sized unit of work
and avoids the specification-implementation divergence that typically occurs
when large formal models are maintained alongside an evolving codebase.

The Agile-SOFL toolchain consists of three components: (i)~a parser that
translates FSF text files into an abstract syntax tree representation suitable
for downstream analysis; (ii)~a language server that provides editing support
(syntax highlighting, completion, and real-time error feedback) through the
Language Server Protocol so that specifications can be authored in standard
editors; and (iii)~a conformance checker that consumes the parsed AST,
generates test inputs via SMT solving, and compares implementation outputs
against the oracle derived from the specification.
\citet{li2023iet} demonstrated that this toolchain detects a substantially higher
proportion of specification faults than manual review alone, particularly for
boundary conditions at scenario guard boundaries.

The survey by \citet{woodcock2009formalsurvey} identifies the primary barrier
to industrial adoption of formal methods as the difficulty of integrating
them with mainstream development workflows.
Agile-SOFL directly attacks this barrier by making formal checking a routine
artefact of sprint-level testing rather than a heavyweight gate at the end of
a waterfall lifecycle.
HSP-Agile builds on this foundation by augmenting the conformance checker
with an LLM-driven candidate synthesis loop, thereby automating not only
\emph{verification} but also \emph{repair} of non-conforming implementations.

% ─────────────────────────────────────────────────────────────────────────────
\section{The Functional Specification Format (FSF)}
\label{sec:bg:fsf}
% ─────────────────────────────────────────────────────────────────────────────

\subsection{Syntax and Structure}

The Functional Specification Format is the primary artifact produced during
the specification phase of an Agile-SOFL sprint \citep{li2022agilesofl, liu2014sofl}.
An FSF document describes a single function $t$ in terms of its \emph{name},
\emph{parameter list}, any \emph{external (state) variables} it reads or writes,
and an ordered sequence of \emph{scenario clauses}.
The general syntactic skeleton is as follows:

\begin{center}
\small\ttfamily
\begin{minipage}{0.82\textwidth}
\obeylines
Function: \textit{name}(\textit{param$_1$}, \ldots, \textit{param$_n$})
~~ext rd \textit{x$_1$}, \textit{x$_2$}, \ldots
~~ext wr \textit{y$_1$}, \textit{y$_2$}, \ldots
~~Scenario S1: \textit{guard$_1$}
~~~~~~\textit{output definitions}
~~Scenario S2: \textit{guard$_2$}
~~~~~~\textit{output definitions}
~~Scenario Sk: others
~~~~~~\textit{default output definitions}
\end{minipage}
\end{center}

\noindent
Each scenario clause $S_i$ consists of (a) a \emph{guard condition}
that is a Boolean expression over the function's parameters and readable
external variables, and (b) an \emph{assignment block} that specifies the
values of output variables and writable external variables when the guard is
satisfied.
The final scenario may use the keyword \texttt{others} as its guard, which
acts as an unconditional catch-all and must appear last.

\subsection{Scenario Ordering Semantics}
\label{sec:bg:fsf:semantics}

The defining semantic feature of FSF is that scenario evaluation proceeds
\emph{top-to-bottom in priority order}: the first scenario whose guard is
satisfied by the current input determines the function's output.
This is analogous to an ordered chain of \texttt{if}--\texttt{else if}
branches rather than a flat set of mutually exclusive conditions.
An input that satisfies both $S_1$'s guard and $S_2$'s guard is handled
exclusively by $S_1$; $S_2$ applies only to inputs in its guard region that
are \emph{not already covered} by higher-priority scenarios.

To make this concrete, consider the following canonical example that appears
throughout this thesis:

\begin{center}
\small\ttfamily
\begin{minipage}{0.85\textwidth}
\obeylines
Function: classify\_risk(a, b, c)
~~Scenario S1: a > 10 and b > 5
~~~~~~risk := "HIGH"; action := "ESCALATE"; score := 90
~~Scenario S2: a > 5
~~~~~~risk := "MEDIUM"; action := "REVIEW"; score := 60
~~Scenario S3: others
~~~~~~risk := "LOW"; action := "MONITOR"; score := 30
\end{minipage}
\end{center}

\noindent
For the input $(a, b, c) = (12, 7, 0)$, both $S_1$'s guard
($12 > 10 \wedge 7 > 5$, i.e.\ true) and $S_2$'s guard ($12 > 5$, i.e.\ true)
are satisfied.
Under FSF semantics the output must be
$\langle\texttt{HIGH}, \texttt{ESCALATE}, 90\rangle$ because $S_1$ has higher priority than $S_2$.
An implementation that tests $a > 5$ before testing $a > 10 \wedge b > 5$
would incorrectly classify this input as \texttt{"MEDIUM"}, even though both
conditions are \emph{individually} satisfiable.
This constitutes a \emph{scenario-order violation}: a category of fault that
is invisible to techniques that treat guard conditions as a flat Boolean
formula without taking their sequential precedence into account.

Figure~\ref{fig:fsf-semantics} (Chapter~\ref{ch:formalization}) illustrates how the three scenarios partition
the input space.
Region $S_1$ (innermost) takes highest priority and is a strict subset of the
region satisfying $S_2$'s guard; the remaining part of $S_2$'s guard region
that lies outside $S_1$ forms the effective domain of $S_2$; and $S_3$ covers
everything else.

The generic three-scenario example in Figure~\ref{fig:fsf-semantics} applies
directly to \code{classify\_risk}: region $S_1$ (red) has highest priority and is a strict subset
of the $S_2$ guard region (orange); $S_3$ (green) covers the remainder.
The precedence relation $S_1 \succ S_2 \succ S_3$ is structural,
not derived from the guards alone.

\subsection{Formal Notation}

We formalise the structure of an FSF task as follows.
A task is a triple $t = (\mathcal{S}_t,\, \mathcal{X}_t,\, \mathcal{Y}_t)$
where:
\begin{itemize}
  \item $\mathcal{X}_t$ is the set of input variable types;
  \item $\mathcal{Y}_t$ is the set of output variable types;
  \item $\mathcal{S}_t = \langle s_1, s_2, \ldots, s_n \rangle$ is the
        \emph{ordered} sequence of scenarios, with $s_1$ having highest priority.
\end{itemize}
Each scenario $s_i = (g_i, \delta_i)$ comprises a guard
$g_i : \mathcal{X}_t \to \{\mathtt{true}, \mathtt{false}\}$ and an output
definition $\delta_i : \mathcal{X}_t \to \mathcal{Y}_t$.
The \emph{oracle} for task $t$ is the function
\[
  g_t(x) \;=\; \delta_j(x)
  \quad\text{where}\quad
  j = \min\bigl\{i \mid g_i(x) = \mathtt{true}\bigr\}.
\]
A conforming implementation $f : \mathcal{X}_t \to \mathcal{Y}_t$ must
satisfy $f(x) = g_t(x)$ for all $x \in \mathcal{X}_t$.
Any input $x^*$ for which $f(x^*) \neq g_t(x^*)$ is a
\emph{counterexample} witnessing non-conformance.

This formulation exposes why scenario ordering matters for conformance
checking: the oracle $g_t$ is \emph{not} equivalent to any unordered Boolean
combination of the guards.
A formal checker that merely asks whether all guards are satisfiable without
encoding priority will accept implementations that violate the ordering
constraint.
The Z3-based conformance checker in HSP-Agile encodes the priority ordering
explicitly, as detailed in Chapter~\ref{ch:formalization}.

It is worth noting the relationship to Dijkstra's guarded commands
\citep{dijkstra1975guarded}.
A guarded command set is \emph{nondeterministic} when multiple guards are
simultaneously true: the system is free to choose any enabled branch.
FSF, by contrast, imposes a deterministic priority ordering that resolves
all overlaps, making the specification suitable for generating a unique
executable oracle.
This distinction is crucial for automated test case generation, which requires
a single ground-truth output for each test input.

% ─────────────────────────────────────────────────────────────────────────────
\section{LLM-Assisted Code Generation}
\label{sec:bg:llm}
% ─────────────────────────────────────────────────────────────────────────────

\subsection{Foundation Models for Code}

The emergence of large language models trained on code corpora has
fundamentally altered the landscape of automated program synthesis.
\citet{chen2021codex} introduced Codex, a GPT-3 variant fine-tuned on
publicly available code from GitHub, and proposed the \emph{pass@k} metric
to quantify the probability that at least one of $k$ independently drawn
samples from the model solves a given programming problem.
Formally, if a model generates $n$ samples of which $c$ pass all unit tests,
the unbiased estimator is:
\[
  \widehat{\mathrm{pass}@k}
  \;=\;
  1 - \frac{\binom{n-c}{k}}{\binom{n}{k}}.
\]
Codex achieved a pass@1 of 28.8\% and pass@100 of 72.3\% on the HumanEval
benchmark, establishing that LLMs can solve a substantial fraction of
introductory programming problems when given multiple attempts.
The \emph{pass@k} framing has since become the standard evaluation
protocol for code generation systems and is adopted, in adapted form, as
one of the baselines in our experimental evaluation.

\citet{openai2023gpt4} documented GPT-4, which substantially improved
on Codex across a range of coding tasks, including competitive programming
problems and formal verification benchmarks.
GPT-4's ability to follow structured prompts and reason about multi-step
logical constraints makes it particularly relevant for the FSF synthesis
task, where the model must respect scenario ordering as well as produce
syntactically correct Python code.

\citet{li2022alphacode} demonstrated that with sufficiently large model scale
and massive code corpora, LLMs can reach competitive-programmer performance on
algorithmic contest problems.
AlphaCode generates large numbers of candidates, filters by a learned
discriminator, and uses clustering to select a diverse final submission set.
While impressive in the competitive programming setting, this approach is
unsuitable for our setting: we require \emph{guaranteed} conformance to a
formal specification, not merely high probability of passing hidden test cases.

\subsection{Evaluation of Deployed Code Generation Tools}

\citet{dakhel2023copilot} conducted a systematic evaluation of GitHub Copilot
on an extensive suite of programming exercises and found that while Copilot
produces correct solutions for straightforward problems, it struggles with
tasks requiring precise boundary handling, multi-condition logic, and adherence
to implicit conventions.
These findings are directly relevant to our setting: FSF scenario guards
involve exactly the kind of boundary and multi-condition logic that Copilot
handles poorly.
In particular, the study observed that Copilot's suggestions were
\emph{semantically plausible}---they looked reasonable to a human reviewer---but
contained subtle logical errors such as off-by-one boundaries and incorrect
operator precedence.
Semantic plausibility without formal correctness is, as we argue in this thesis,
the core challenge of LLM-based formal specification implementation.

\subsection{Limitations of LLMs for Formal Specification Compliance}

Three systematic failure modes arise when LLMs generate implementations of
FSF specifications without additional guidance \citep{dakhel2023copilot, lynn2024verifierloop}:

\begin{enumerate}
  \item \textbf{Scenario-order blindness.}
        LLMs trained on code learn the surface syntax of
        \texttt{if}--\texttt{elif} chains but do not learn to interpret
        the \emph{specification-level} priority ordering of FSF scenarios.
        When two scenarios have overlapping guards, the model's training
        distribution biases it towards the condition that appears more
        ``natural'' in code rather than the one mandated by the specification.

  \item \textbf{Boundary hallucination.}
        Numeric comparison operators (\texttt{>}, \texttt{>=}, \texttt{==})
        are frequently confused at specification boundaries.
        For example, a specification with guard $a > 10$ may be implemented
        with $a \geq 10$, introducing a one-point discrepancy that is
        invisible to random testing but detectable by an SMT-generated boundary
        test case.

  \item \textbf{Spurious completeness.}
        LLMs sometimes generate implementations that always return a
        hard-coded success value for some output fields, producing code that
        passes superficial tests but violates the specification for inputs
        that should trigger non-default behaviour.
        This pattern corresponds to our high-severity requirement fault RF07
        (\emph{unconditional success}) discussed in Section~\ref{sec:bg:patterns}.
\end{enumerate}

\noindent
The standard approach to improving LLM code quality is to provide test-case
feedback in a repair loop \citep{chen2021codex}.
In this paradigm, a set of unit tests is executed against each generated
candidate; failing test cases are included in the prompt for a subsequent
generation attempt.
This is the strategy employed by our B2 (test-feedback) baseline.
However, the test-feedback loop has a fundamental limitation in the FSF
context: the tests themselves must be \emph{derived from} the specification,
and if they fail to cover scenario boundary inputs, the loop may converge on
a candidate that passes all available tests while still violating the
specification on uncovered inputs.
HSP-Agile addresses this by using Z3 to \emph{synthesise} boundary-hitting
test cases that are guaranteed to expose scenario-order violations, rather
than relying on a fixed test suite.

% ─────────────────────────────────────────────────────────────────────────────
\section{Counterexample-Guided Synthesis (CEGIS)}
\label{sec:bg:cegis}
% ─────────────────────────────────────────────────────────────────────────────

\subsection{The Classical CEGIS Framework}

Counterexample-Guided Inductive Synthesis (CEGIS) was introduced as a
principled framework for program synthesis by \citet{solarlezama2008sketching}.
The core idea is to decompose the synthesis problem into two interacting
sub-problems: a \emph{synthesis} sub-problem that produces a candidate
program satisfying a partial specification, and a \emph{verification}
sub-problem that checks whether the candidate satisfies the full specification
and, if not, produces a \emph{counterexample}---a concrete input that
witnesses the violation.
The counterexample is fed back to the synthesiser as an additional constraint,
and the loop repeats until either a correct program is found or a
resource budget is exhausted.

In the original formulation of \citet{solarlezama2008sketching}, both the
synthesiser and the verifier are SAT/SMT solvers: the synthesiser fills in
``holes'' in a sketch (a program template with unknown sub-expressions) while
the verifier checks the completed program against a universal quantification
over inputs.
The separation between synthesiser and verifier is essential because it
allows each component to be specialised: the synthesiser operates over
a \emph{finite} set of counterexamples and can use polynomial-time
procedures, while the verifier handles the universally quantified
correctness condition.

\subsection{LLM-Augmented CEGIS Variants}

The classical CEGIS framework assumes that the synthesiser operates over a
fixed language of program templates.
With the advent of code-generating LLMs, it became natural to replace the
template-based synthesiser with a language model that can generate
free-form code from a natural-language description.
This hybrid approach retains the formal verifier component, which provides
the soundness guarantee, while delegating the creative search over the
program space to the LLM, which is better at generating plausible code
from high-level descriptions.

\citet{loughman2024dafnybench} introduced DafnyBench, a benchmark for
evaluating LLM-based synthesis of verified Dafny programs.
Their evaluation shows that state-of-the-art LLMs can synthesise many
simple Dafny proofs but consistently fail on programs requiring intricate
loop invariants or complex pre-/post-condition reasoning, precisely the
cases where the formal verifier provides the most value.

\citet{lynn2024verifierloop} developed a ``verifier-in-the-loop'' methodology
for TOSEM in which an LLM generates code and a formal tool checks correctness,
with failures reported as structured feedback.
Their system demonstrated improved correctness rates over baseline LLM
generation across several specification languages, but did not address
the scenario-ordering semantics of FSF nor provide the pattern-guard
layer that blocks structurally pathological candidates before formal
checking.

\subsection{HSP-Agile's CEGIS Adaptation}

HSP-Agile instantiates the CEGIS loop with three specialised components
tailored to the FSF setting.
The \emph{synthesiser} is an LLM (GPT-4) prompted with the FSF specification
and the accumulated repair context from previous iterations.
The \emph{verifier} is the Z3-based formal conformance checker, which executes
the candidate on the SMT-generated test suite $\mathcal{D}(t)$ and identifies
failing inputs.
The \emph{feedback aggregator} translates failing test cases into a structured
natural-language repair context that includes the counterexample input, the
expected output $g_t(x^*)$, and the observed output $f_{c_k}(x^*)$, together
with any high-severity requirement patterns detected in the candidate's source
code.

Figure~\ref{fig:cegis-loop} illustrates the four-component loop.
The key departure from prior LLM+verifier approaches is that the feedback
aggregator communicates \emph{both} formal counterexamples \emph{and}
pattern-guard violations, enabling the LLM to simultaneously correct
logical errors and structural anti-patterns in a single repair step.

\begin{figure}[t]
  \centering
  \tikzinput{diagrams/tikz/cegis_loop}
  \caption{The HSP-Agile CEGIS loop.
           The LLM synthesiser produces a candidate $c_k$ from the FSF
           specification $\mathcal{S}_t$ and accumulated repair context.
           The formal conformance checker verifies $c_k$ against the
           Z3-generated test suite; failing cases are extracted into
           counterexamples $(x^*, g_t(x^*), f_{c_k}(x^*))$.
           The feedback aggregator packages counterexamples and detected
           anti-patterns into the repair context for the next iteration.
           If all tests pass and pattern guards are clear, $c_k$ is accepted.}
  \label{fig:cegis-loop}
\end{figure}

\noindent
The budget parameter $K$ controls the maximum number of repair iterations.
In our experiments we set $K = 5$; if no conforming candidate is found within
$K$ attempts, the system returns the candidate with the highest number of
passing tests as a best-effort result.
This is in line with the budget-bounded CEGIS strategy discussed
by \citet{barrett2018smt} in the context of bounded model checking.

% ─────────────────────────────────────────────────────────────────────────────
\section{SMT Solving and Formal Verification}
\label{sec:bg:smt}
% ─────────────────────────────────────────────────────────────────────────────

\subsection{Satisfiability Modulo Theories}

Satisfiability Modulo Theories (SMT) extends propositional satisfiability
(SAT) by allowing formulas to contain atoms drawn from background
\emph{theories} such as linear arithmetic over integers and reals,
array theory, uninterpreted functions, and bitvectors
\citep{barrett2018smt}.
An SMT solver takes as input a formula $\varphi$ in a combination of
supported theories and either returns a \emph{satisfying assignment} (a
concrete valuation of all variables that makes $\varphi$ true) or proves
that $\varphi$ is \emph{unsatisfiable}.
The satisfying assignment is directly usable as a concrete test input,
which is the mechanism exploited by HSP-Agile's test case generation.

\citet{demoura2008z3} introduced Z3, which has become the de facto standard
SMT solver in software verification research.
Z3 supports the theories of linear integer and real arithmetic, arrays,
datatypes, and quantifiers, and exposes both a command-line interface and
programmatic APIs in Python, C, and other languages.
Its incremental solving capabilities allow a verifier to assert and retract
constraints during a single session, an important efficiency property for
HSP-Agile's conformance checker, which processes each scenario's guard
incrementally during test case synthesis.

\subsection{Test Case Generation via SMT}

Given an FSF task $t = (\langle s_1, \ldots, s_n \rangle, \mathcal{X}_t,
\mathcal{Y}_t)$, the conformance checker constructs the test suite
$\mathcal{D}(t)$ by solving, for each scenario $s_i$, the formula:
\[
  \varphi_i \;=\; g_i(x) \;\wedge\; \bigwedge_{j < i} \neg g_j(x).
\]
The formula $\varphi_i$ characterises the \emph{effective} guard of $s_i$:
inputs that satisfy $s_i$'s guard but are \emph{not} already covered by any
higher-priority scenario.
A satisfying assignment $x_i^*$ to $\varphi_i$ is a concrete test input that,
under the oracle, must produce output $\delta_i(x_i^*)$, i.e.\ the output
defined by scenario $s_i$.
Testing a candidate $f$ against $x_i^*$ and comparing $f(x_i^*)$ with
$g_t(x_i^*) = \delta_i(x_i^*)$ directly checks whether $f$ respects the
priority ordering of the specification.

This approach is related to, but distinct from, property-based testing
frameworks such as QuickCheck \citep{claessen2000quickcheck}.
QuickCheck generates random inputs from a user-specified distribution and
checks them against a user-defined property.
While highly effective for detecting coarse logical errors, random generation
does not guarantee coverage of scenario-boundary inputs and cannot encode the
FSF priority ordering as a test generation strategy.
SMT-based generation, by contrast, is \emph{targeted}: it produces inputs
that are guaranteed to lie in the effective guard region of each scenario,
ensuring that the boundary between adjacent scenarios is always represented
in the test suite regardless of the distribution of random inputs.

\subsection{Relationship to Full Formal Verification}

The Dafny language \citep{leino2010dafny} supports full deductive verification
via annotated programs with pre- and post-conditions, loop invariants, and
termination arguments.
A Dafny program that passes verification is \emph{proved correct} with respect
to its specification under all possible inputs, not just the finite set of
inputs in a test suite.
This provides stronger guarantees than the bounded checking strategy of
HSP-Agile.

However, Dafny verification requires the developer to write and maintain
annotations (loop invariants, in particular) that can be as complex as the
program itself.
For the Agile-SOFL context---where specifications are written by engineers
with limited formal methods background and where sprint velocity is a key
concern---the annotation burden of full verification is prohibitive.
HSP-Agile deliberately occupies a middle ground: stronger than
random testing because test inputs are SMT-generated to target scenario
boundaries, but lighter than full verification because no annotations are
required from the developer.
The scope of the prevention guarantee is bounded by the coverage of the
generated test suite, which is a known and accepted limitation discussed
further in Chapter~\ref{ch:discussion}.

% ─────────────────────────────────────────────────────────────────────────────
\section{Requirement-Level Error Pattern Detection}
\label{sec:bg:patterns}
% ─────────────────────────────────────────────────────────────────────────────

\subsection{Static Analysis and Bug Pattern Catalogues}

Static analysis tools detect potential defects in source code without
executing the program by analysing its syntactic and structural properties.
\citet{hovemeyer2004findbugs} introduced FindBugs, one of the most widely
deployed static analysis tools for Java, which encodes a catalogue of
\emph{bug patterns}---recurring structural templates that correlate strongly
with runtime defects.
FindBugs analyses compiled bytecode using a visitor-based abstract
interpretation \citep{cousot1977abstract} and reports violations of patterns
such as null pointer dereferences, incorrect \texttt{equals}
implementations, and misuse of floating-point comparisons.

\citet{ayewah2008findbugs} evaluated FindBugs on production code from
Sun Microsystems and Google and found that a substantial fraction of its
warnings corresponded to genuine defects with real-world impact, though the
false-positive rate varied considerably across pattern categories.
This observation motivates a tiered severity model for pattern warnings:
high-severity patterns (low false-positive rate, high defect correlation)
should block acceptance, while lower-severity patterns (higher false-positive
rate) should be reported as warnings only.

\subsection{Requirement-Level Patterns in FSF Implementations}

The key conceptual contribution of HSP-Agile's pattern guard layer is the
recognition that the relevant bug patterns for FSF implementations are
\emph{requirements-level violations}, not merely code smells.
A code smell such as an overly long method or a deeply nested conditional is
a maintainability concern; it does not necessarily indicate a logical error.
A requirements-level pattern violation, by contrast, signals that the
implementation cannot possibly conform to the FSF specification regardless
of the specific inputs considered.

The HSP-Agile pattern catalogue RF01--RF14 was assembled by analysing the
failure modes observed during LLM-based FSF implementation \citep{li2023iet, hovemeyer2004findbugs}.
Representative high-severity patterns include RF01 (missing precondition check),
RF02 (wrong default branch), RF05 (swapped comparison), and RF07 (unconditional
success); the full catalogue is documented in Appendix~\ref{app:patterns}.

\noindent
Patterns RF01--RF14 are detected by static analysis of the candidate AST
\citep{cousot1977abstract}.
Each pattern checker is a visitor that reports a violation when the
structural signature matches.
High-severity patterns ($\mathcal{P}^H$) block acceptance and are included
in the CEGIS repair context even when all formal test cases pass,
covering input regions the finite test suite may omit.

% ─────────────────────────────────────────────────────────────────────────────
\section{Mutation Testing for Specification Assessment}
\label{sec:bg:mutation}
% ─────────────────────────────────────────────────────────────────────────────

\subsection{Foundations of Mutation Testing}

Mutation testing was introduced by \citet{demillo1978mutation} as a method
for assessing the adequacy of a test suite.
The central idea is to inject small, well-defined syntactic changes
(\emph{mutations}) into a program to produce a set of \emph{mutants}, and
then measure what fraction of mutants are \emph{killed} by the test suite
(i.e.\ cause at least one test to fail).
A mutant that is not killed by any test is called a \emph{surviving mutant};
a high rate of survivors suggests that the test suite lacks inputs
that can distinguish the correct program from plausibly incorrect variants.
\citet{offutt1992mutation} formalised a set of \emph{mutation operators}---
templates for syntactic modifications---and showed that killing a carefully
chosen set of \emph{sufficient operators} implies that the test suite also
kills a much larger class of higher-order mutants.

\citet{ammann2016testing} provides a comprehensive treatment of mutation
testing, establishing the relationship between mutation coverage and
various logical coverage criteria.
The central result is that achieving full mutation coverage is strictly
stronger than achieving full condition/decision coverage (MC/DC), making
mutation testing a useful complement to structural coverage analysis.

\subsection{Specification and Implementation Mutation Operators}

HSP-Agile uses mutation testing at two levels that correspond to the two
components being assessed.

\paragraph{Specification-level mutation} evaluates whether the FSF specification
itself is sufficiently discriminating: does the Z3-generated test suite
$\mathcal{D}(t)$ kill mutants of the oracle $g_t$?
The specification mutation operators used in our experiments are:

\begin{itemize}
  \item \textbf{SNO} (Scenario Negate Once): negate the guard condition of
        one scenario, e.g.\ $a > 10 \wedge b > 5$ becomes
        $\neg(a > 10 \wedge b > 5)$.
  \item \textbf{ORO} (Operator Replacement Ordinal): replace a relational
        operator in a guard, e.g.\ $>$ becomes $\geq$, $<$, or $=$.
  \item \textbf{SPO} (Scenario Priority Order): swap the priority of two
        adjacent scenarios, e.g.\ move $S_2$'s assignment block to $S_1$
        and vice versa.
\end{itemize}

A high specification-mutation kill rate indicates that the test suite
covers the boundaries between scenarios and is sensitive to priority-order
violations---precisely the class of error that HSP-Agile is designed to
prevent.

\paragraph{Implementation-level mutation} evaluates whether the generated
implementation $c_k$ is robust to the kinds of implementation errors that
LLMs are prone to make.
The implementation mutation operators include:

\begin{itemize}
  \item \textbf{ICO} (Invert Condition Once): invert a Boolean sub-expression
        in a conditional guard, e.g.\ \texttt{a > 10} becomes
        \texttt{not (a > 10)}.
  \item \textbf{WRO} (Wrong Return Output): replace the return value in one
        scenario branch with a value from a different scenario's output block.
  \item \textbf{BCO} (Boundary Change Off-by-one): replace a strict
        comparison operator with a non-strict one or vice versa at a scenario
        boundary.
\end{itemize}

\subsection{Prevention-Oriented Metrics}

Standard mutation testing metrics measure test suite quality.
In the HSP-Agile context, mutation testing is repurposed as a
\emph{prevention assessment} tool: we measure not how well a test suite
detects faults in an existing program, but how well the pipeline \emph{prevents}
non-conforming implementations from being accepted.

Two prevention metrics are central to our evaluation:

\begin{description}
  \item[Prevention Detection Rate (PDR).]
        The fraction of deliberately introduced faults (instantiated
        via implementation mutations) that the pipeline correctly rejects.
        A PDR of 1.0 means the pipeline rejects every non-conforming
        candidate; a lower PDR indicates scenarios where the pipeline
        admits faulty code.

  \item[False Accept Rate (FAR).]
        The fraction of correct implementations (true conforming candidates)
        that the pipeline incorrectly rejects.
        A low FAR is essential for practical usability: if the pipeline
        rejects too many correct implementations, it creates excessive repair
        overhead and undermines developer trust.
\end{description}

\noindent
These metrics mirror the sensitivity--specificity trade-off in
binary classification and are discussed further in
Chapter~\ref{ch:formalization}.
\citet{okun2007mutation} used analogous metrics (detection rate and
false alarm rate) in the context of static analysis tool evaluation,
and their analysis motivates our choice of PDR and FAR as the
primary prevention-quality measures.

% ─────────────────────────────────────────────────────────────────────────────
\section{Related Work and Research Gap}
\label{sec:bg:related}
% ─────────────────────────────────────────────────────────────────────────────

\subsection{Formal Methods in Software Engineering Practice}

The integration of formal methods into mainstream software development has
been a persistent research challenge for four decades.
\citet{woodcock2009formalsurvey} conducted an extensive survey of industrial
formal method applications and identified recurring themes: the most
successful deployments combine formal specification with automated analysis
tools (model checkers, SMT solvers) and integrate into existing development
workflows rather than replacing them.
Projects that adopted formal methods in a narrow, high-value layer
of a larger system---rather than demanding that all code be formally verified
---achieved the best balance of rigour and productivity.
Agile-SOFL is consistent with this prescription; HSP-Agile takes it further
by automating the synthesis and repair loop so that formal checking requires
no manual annotation effort.

Requirements engineering research \citep{pohl2010requirements} emphasises
the importance of detecting ambiguities and inconsistencies at the
specification level rather than discovering them through integration testing.
Specification-level error detection is substantially cheaper: the cost of
correcting a requirements defect grows superlinearly with how late in the
development lifecycle it is found.
HSP-Agile's pattern guard layer operates at the requirements level,
checking whether the generated implementation's structure is consistent with
the FSF specification before even executing a single test case.

\subsection{Approaches to LLM-Verified Code Synthesis}

Several recent systems have combined LLM-based code generation with
formal or semi-formal verification.
We survey the most closely related approaches and compare them
systematically in Table~\ref{tab:related-work}.

\paragraph{DafnyBench \citep{loughman2024dafnybench}.}
DafnyBench evaluates LLMs on the task of synthesising Dafny programs with
correct verification annotations.
The Dafny verifier provides full deductive proofs of functional correctness,
offering the strongest possible guarantee.
However, DafnyBench's repair loop relies on the Dafny compiler's error messages
rather than semantic counterexamples: when verification fails, the error
output is appended to the prompt but does not identify a concrete
failing input that illustrates the logical discrepancy.
Moreover, DafnyBench has no pattern guard layer: structurally pathological
candidates (e.g.\ those with missing output fields) are submitted to the
full verifier without a cheap first-pass filter.
DafnyBench's evaluation metric is pass@k, which does not distinguish between
approaches that achieve correctness through rigorous formal reasoning and
those that do so by statistical luck.

\paragraph{Lynn et al.\ TOSEM \citep{lynn2024verifierloop}.}
The verifier-in-the-loop methodology of Lynn et al.\ is the closest prior
work to HSP-Agile.
Their system uses a formal tool (a model checker or type checker, depending
on the specification language) as an oracle and feeds failure reports back to
the LLM.
The key differences from HSP-Agile are: (i)~their target specification language
is not FSF and does not encode scenario priority ordering; (ii)~the repair loop
does not include a pattern guard layer; and (iii)~their evaluation metric is
a correctness rate over a fixed benchmark, which does not distinguish
false-accept and false-reject errors in the way that PDR and FAR do.
Their reported improvement over non-feedback baselines is substantial (up to
23 percentage points) but their benchmark does not include tasks with
overlapping guard conditions, where scenario-order violations are most likely.

\paragraph{Standard LLM baseline (B1).}
The simplest baseline is to prompt a code-generating LLM with the FSF
specification and accept the first generated candidate without any
verification or feedback.
This approach has no formal specification in the conformance sense:
the LLM is asked to follow the specification as a natural-language prompt
but there is no formal check that it does so.
The pass rate on our benchmark for this baseline quantifies how often LLMs
spontaneously produce conforming implementations, providing a measure of the
baseline capability that HSP-Agile must improve upon.

\paragraph{Test-feedback baseline (B2).}
The test-feedback baseline adds an iterative repair loop to B1: a fixed
set of unit tests (not SMT-generated) is executed against each candidate,
and failing tests are appended to the repair prompt.
This baseline is representative of widely deployed LLM coding assistants that
incorporate automated test execution feedback.
It does not use a formal specification for test generation and has no pattern
guard layer.
The comparison between B2 and HSP-Agile isolates the contribution of
SMT-based test case generation and pattern guards over purely test-driven
feedback.

\subsection{Structured Comparison}

\begin{table}[t]
  \centering
  \caption{Structured comparison of related approaches to LLM-assisted
           specification-conformant code generation.
           Checkmarks (\checkmark) indicate the feature is fully present;
           dashes (--) indicate it is absent;
           ``partial'' indicates limited support.
           Column abbreviations:
           F.\,Spec = formal specification;
           Repair = repair loop;
           Guard = pattern guard;
           Prev.\,Metr.\ = prevention metrics.}
  \label{tab:related-work}
  \thesistable{%
    \footnotesize
    \renewcommand{\arraystretch}{1.12}
    \begin{tabular*}{\linewidth}{@{\extracolsep{\fill}}l c c c c c@{}}
      \toprule
      \textbf{Approach}
        & \textbf{F.\,Spec}
        & \textbf{LLM}
        & \textbf{Repair}
        & \textbf{Guard}
        & \textbf{Prev.\,Metr.} \\
      \midrule
      DafnyBench \citep{loughman2024dafnybench}
        & \checkmark\ (Dafny)
        & \checkmark
        & partial
        & --
        & pass@k \\
      Lynn et al.\ TOSEM \citep{lynn2024verifierloop}
        & \checkmark
        & \checkmark
        & partial
        & --
        & correctness \\
      Standard LLM (B1)
        & --
        & \checkmark
        & --
        & --
        & pass rate \\
      Test-feedback (B2)
        & --
        & \checkmark
        & \checkmark\ (tests)
        & --
        & pass rate \\
      \midrule
      \textbf{HSP-Agile (ours)}
        & \checkmark\ \textbf{(FSF)}
        & \checkmark
        & \checkmark\ \textbf{(CEGIS)}
        & \checkmark
        & \textbf{PDR, FAR, Conf} \\
      \bottomrule
    \end{tabular*}%
  }
\end{table}

\noindent
Table~\ref{tab:related-work} reveals that no prior approach combines all four
dimensions simultaneously: formal FSF specification with scenario ordering,
LLM-driven synthesis, a CEGIS repair loop with SMT-generated counterexamples,
and a pattern guard layer.
Prior work evaluates quality primarily through pass@k or aggregate correctness
rates \citep{chen2021codex, loughman2024dafnybench}, neither of which
distinguishes false acceptances from false rejections.
HSP-Agile introduces PDR and FAR as prevention-quality metrics that measure
both error types directly.

\subsection{Research Gap}

Synthesising the observations above, the research gap addressed by this
thesis can be stated precisely:

\begin{enumerate}
  \item \textbf{Formal scenario semantics.}
        Existing LLM-based code generation tools---including those with
        formal verification components---do not account for FSF scenario
        priority ordering.
        Tests generated from flat guard conditions or from natural-language
        descriptions will systematically miss the class of scenario-order
        violations described in Section~\ref{sec:bg:fsf:semantics}.

  \item \textbf{Requirement-level prevention.}
        Static analysis tools (FindBugs and its successors) operate at
        the code level and are not specification-aware.
        No prior tool provides a pattern guard layer whose patterns are
        derived from the structure of a formal FSF specification and
        evaluated before formal test execution.

  \item \textbf{Prevention-oriented evaluation.}
        The LLM code generation literature evaluates models using
        generation quality metrics (pass@k, BLEU) that do not address
        the binary question relevant to formal specification contexts:
        ``Does this candidate conform to the specification, and at what
        cost in false rejections?''
        PDR and FAR provide exactly this distinction.

  \item \textbf{Agile integration.}
        Prior formal verification tools for LLM-generated code
        \citep{loughman2024dafnybench, lynn2024verifierloop} are evaluated on
        academic benchmarks; none reports integration with sprint-based
        workflows \citep{beck2001xp, li2022agilesofl}.
        HSP-Agile is designed for Agile-SOFL sprints, and our benchmark
        comprises sprint-sized FSF tasks.
\end{enumerate}

\noindent
HSP-Agile addresses all four gaps through a unified pipeline whose components
are described formally in Chapter~\ref{ch:formalization} and implemented
and evaluated in Chapters~\ref{ch:implementation}--\ref{ch:experiments}.
The following chapter establishes the mathematical foundations that underpin
each component of the pipeline.
